% Standalone problem document, generated from the adjacent README.md.
% Compile with XeLaTeX. Mathematical content is maintained in Markdown.
\documentclass[11pt,a4paper]{article}
\usepackage[fontsize=10.5pt]{scrextend}
\usepackage[margin=22mm,headheight=15pt,headsep=9mm,footskip=11mm]{geometry}
\usepackage{fontspec}
\setmainfont{texgyrepagella}[Extension=.otf,UprightFont=*-regular,BoldFont=*-bold,ItalicFont=*-italic,BoldItalicFont=*-bolditalic]
\setsansfont{texgyreheros}[Extension=.otf,UprightFont=*-regular,BoldFont=*-bold,ItalicFont=*-italic,BoldItalicFont=*-bolditalic]
\setmonofont{texgyrecursor}[Extension=.otf,UprightFont=*-regular,BoldFont=*-bold,ItalicFont=*-italic,BoldItalicFont=*-bolditalic,Scale=MatchLowercase]
\usepackage{amsmath,amssymb}
\usepackage{unicode-math}
\setmathfont{texgyrepagella-math.otf}
\usepackage{xcolor,microtype,enumitem,needspace,fancyhdr}
\usepackage{bookmark}
\definecolor{ink}{HTML}{17344B}
\definecolor{accent}{HTML}{197D80}
\definecolor{muted}{HTML}{596773}
\hypersetup{colorlinks=true,linkcolor=accent,urlcolor=accent,pdftitle={NR-02 - Additivity
for Cartesian-product slack
matrices},pdfauthor={Open Problems in Numerical Linear Algebra}}
\urlstyle{same}
\setlength{\parindent}{0pt}
\setlength{\parskip}{5pt plus 1pt minus 1pt}
\linespread{1.04}
\setlength{\emergencystretch}{3em}
\widowpenalty=10000
\clubpenalty=10000
\displaywidowpenalty=10000
\allowdisplaybreaks[1]
\setlist{nosep,leftmargin=1.5em}
\providecommand{\tightlist}{\setlength{\itemsep}{0pt}\setlength{\parskip}{0pt}}
\providecommand{\pandocbounded}[1]{#1}
\pagestyle{fancy}
\fancyhf{}
\fancyhead[L]{\sffamily\footnotesize\color{muted}OPEN PROBLEMS IN NUMERICAL LINEAR ALGEBRA}
\fancyhead[R]{\sffamily\bfseries\footnotesize\color{accent}NR-02}
\fancyfoot[L]{\parbox[b]{0.9\textwidth}{\sffamily\fontsize{8}{10}\selectfont\color{muted}Editorial ratings; a bounded search cannot certify that a problem remains unsolved.\\Literature check: 2026-09-10}}
\fancyfoot[R]{\sffamily\footnotesize\color{muted}\thepage}
\renewcommand{\headrulewidth}{0.3pt}
\renewcommand{\headrule}{\hbox to\headwidth{\color{accent}\leaders\hrule height \headrulewidth\hfill}}
\makeatletter
\renewcommand{\section}{\Needspace{5\baselineskip}\@startsection{section}{1}{0pt}{12pt plus 2pt minus 2pt}{4pt}{\sffamily\bfseries\large\color{ink}}}
\renewcommand{\subsection}{\Needspace{4\baselineskip}\@startsection{subsection}{2}{0pt}{9pt plus 1pt minus 1pt}{3pt}{\sffamily\bfseries\normalsize\color{ink}}}
\makeatother
\setcounter{secnumdepth}{0}
\begin{document}
{\sffamily\small\color{muted}Nonnegative and positive
factorizations\par}
\vspace{3pt}
{\raggedright\hyphenpenalty=10000\exhyphenpenalty=10000\sffamily\bfseries\fontsize{21}{24}\selectfont\color{ink}Additivity
for Cartesian-product slack matrices\par}
\vspace{7pt}
{\sffamily\small\textbf{Difficulty:} extreme\quad\textbf{Importance:} interesting
to the community\par}
{\sffamily\small\bfseries\color{accent}Status: Partially resolved\par}
\vspace{4pt}
{\color{accent}\hrule height 0.8pt}
\vspace{6pt}
\textbf{Rating rationale:} Extreme because additivity across arbitrary
polytopes is a general structural barrier for extension complexity;
community importance links optimal factorizations with compositional
linear optimization models.

\textbf{Area:} nonnegative rank and polyhedral optimization

\subsection{Context and notation}\label{context-and-notation}

All factorizations are over the real numbers. For
\(X\in\mathbb R_{\ge0}^{m\times n}\), define

\[
\operatorname{rank}_+(X)=\min\{r\ge0:X=WH,\quad
W\in\mathbb R_{\ge0}^{m\times r},\ H\in\mathbb R_{\ge0}^{r\times n}\}.
\]

\subsection{Problem statement}\label{problem-statement}

Let \(P\subset\mathbb R^d\) and \(Q\subset\mathbb R^e\) be
full-dimensional polytopes with \(d,e\ge1\). Let
\(S\in\mathbb R_{\ge0}^{m\times n}\) and
\(T\in\mathbb R_{\ge0}^{p\times q}\) be their slack matrices: rows
correspond to all facets in irredundant inequality descriptions, columns
to all vertices, and each entry is the right-hand side minus the
left-hand side of that facet inequality at that vertex. Construct
\(C\in\mathbb R_{\ge0}^{(m+p)\times nq}\) by giving it one column for
each pair \((i,j)\) of vertices:

\[
C[:,(i,j)]=\begin{pmatrix}S[:,i]\\T[:,j]\end{pmatrix}.
\]

\subsubsection{Question}\label{question}

Is

\[
\operatorname{rank}_+(C)=\operatorname{rank}_+(S)+\operatorname{rank}_+(T)
\]

always true? The hypotheses that \(S,T\) are polytope slack matrices are
part of the question. Through the slack-factorization theorem, this also
asks whether the minimum number of inequalities in an extended
formulation is additive under Cartesian products.

\subsection{References}\label{references}

Gillis,
\href{https://orbi.umons.ac.be/bitstream/20.500.12907/42337/1/NMFbook_SIAM_reprint.pdf}{\emph{Nonnegative
Matrix Factorization}}, §3.6.5, pp.~93--94. Hans Raj Tiwary, Stefan
Weltge, and Rico Zenklusen,
\href{https://arxiv.org/abs/1702.01959}{\emph{Extension complexities of
Cartesian products involving a pyramid}}, Information Processing Letters
\textbf{128} (2017), 11--13, introduction and main theorem.

\subsection{Status check --- 2026-09-10}\label{status-check-2026-09-10}

Rechecked
\href{https://arxiv.org/abs/1702.01959}{Tiwary--Weltge--Zenklusen} and
\href{https://www.cargese.org/2022/open-problems.pdf}{Weltge's 2022
workshop question}, then searched for later product-additivity proofs
and counterexamples. Equality is proved whenever one factor is a
pyramid, a substantive portion of the displayed target. The unrestricted
equality remains open in the latest explicit source located; no
subsequent full resolution was found. The 2022 date limits the strength
of the current-status evidence.
\end{document}
